\documentclass[letterpaper,10pt,english]{article}
\usepackage{%
amsfonts,%
amsmath,%
amssymb,%
amsthm,%
babel,%
bbm,%
%biblatex,%
caption,%
centernot,%
color,%
enumerate,%
%enumitem,%
epsfig,%
epstopdf,%
etex,%
fancybox,%
framed,%
fullpage,%
%geometry,%
graphicx,%
hyperref,%
latexsym,%
mathptmx,%
mathtools,%
multicol,%
paralist,%
pgf,%
pgfplots,%
pgfplotstable,%
pgfpages,%
proof,%
psfrag,%
%subfigure,%
tcolorbox,%
tfrupee,%
tikz,%
times,%
%ulem,%
url,%
xcolor,%
mathpazo,%
}

\definecolor{shadecolor}{gray}{.95}%{rgb}{1,0,0}
\usepackage[margin=1in,top=0.75in]{geometry}
\usepackage[mathscr]{eucal}
\usepgflibrary{shapes}
\usepgfplotslibrary{fillbetween}
\usetikzlibrary{%
arrows,%
backgrounds,%
chains,%
decorations.pathmorphing,% /pgf/decoration/random steps | erste Graphik
decorations.text,% 
matrix,%
positioning,% wg. " of "
fit,%
patterns,%
petri,%
plotmarks,%
scopes,%
shadows,%
shapes.misc,% wg. rounded rectangle
shapes.arrows,%
shapes.callouts,%
shapes%
}

%\pgfplotsset{compat=newest} %<------ Here
\pgfplotsset{compat=1.11} %<------ Or use this one

\theoremstyle{plain}
\newtheorem{thm}{Theorem}[section]
\newtheorem{lem}[thm]{Lemma}
\newtheorem{prop}[thm]{Proposition}
\newtheorem{cor}[thm]{Corollary}
\newtheorem{clm}[thm]{Claim}

\theoremstyle{definition}
\newtheorem{defn}[thm]{Definition}
\newtheorem{conj}[thm]{Conjecture}
\newtheorem{exmp}[thm]{Example}
\newtheorem{axiom}[thm]{Axiom}
\newtheorem{assum}[thm]{Assumption}
\newtheorem{exerc}[thm]{Exercise}
\newtheorem*{defin}{Definition}

\theoremstyle{remark}
\newtheorem{rem}{Remark}
\newtheorem{note}{Note}

\newcommand{\iid}{\emph{i.i.d.}\ }
\newcommand{\Cov}{\operatorname{Cov}}
%\newcommand{\det}{\operatorname{det}}
%\newcommand{\Real}{\mathbb{R}}
\newcommand{\tr}{\operatorname{tr}}
\newcommand{\Var}{\operatorname{Var}}
\newcommand{\cov}{\operatorname{cov}}

%\renewcommand{\proof}[1]{\begin{proof}#1\end{proof}}
\newcommand{\EQ}[1]{\begin{equation*}#1\end{equation*}}
\newcommand{\EQN}[1]{\begin{equation}#1\end{equation}}
\newcommand{\eq}[1]{\begin{align*}#1\end{align*}}
\newcommand{\meq}[2]{\begin{xalignat*}{#1}#2\end{xalignat*}}
\newcommand{\norm}[1]{\left\lVert#1\right\rVert}
\newcommand{\inner}[1]{\left\langle #1\right\rangle}
\newcommand{\abs}[1]{\left\lvert#1\right\rvert}
\newcommand{\expect}[1]{\mathbb{E}\left[{#1}\right]}
\newcommand{\prob}[1]{\mathbb{P}\left[{#1}\right]}
\newcommand{\given}{\; \big\vert \;} 
\newcommand{\set}[1]{\left\{#1\right\}} 
\newcommand{\indicator}[1]{1_{\set{#1}}} 
\newcommand{\Ind}[1]{\mathbbm{1}_{#1}} 
\newcommand{\SetIn}[1]{\mathbbm{1}_{\set{#1}}} 
\newcommand{\red}[1]{\textcolor{red}{#1}} 
\newcommand{\floor}[1]{\left\lfloor#1\right\rfloor}
\newcommand{\ceil}[1]{\left\lceil#1\right\rceil}


\newcommand{\C}{\mathbb{C}}
\newcommand{\D}{\mathbb{D}}
\newcommand{\E}{\mathbb{E}}
\newcommand{\F}{\mathbb{F}}
\newcommand{\N}{\mathbb{N}}
\renewcommand{\P}{\mathbb{P}}
\newcommand{\Q}{\mathbb{Q}}
\newcommand{\R}{\mathbb{R}}
\newcommand{\Z}{\mathbb{Z}}

\newcommand{\bA}{\mathbf{A}}
\newcommand{\bI}{\mathbf{I}}
\newcommand{\bU}{\mathbf{1}}
\newcommand{\bx}{\mathbf{x}}
\newcommand{\bX}{\mathbf{X}}
\newcommand{\bZ}{\mathbf{Z}}


\newcommand{\cA}{\mathcal{A}}
\newcommand{\cB}{\mathcal{B}}
\newcommand{\cC}{\mathcal{C}}
\newcommand{\cF}{\mathcal{F}}
\newcommand{\cG}{\mathcal{G}}
\newcommand{\cM}{\mathcal{M}}
\newcommand{\cN}{\mathcal{N}}
\newcommand{\cP}{\mathcal{P}}
\newcommand{\cT}{\mathcal{T}}
\newcommand{\cX}{\mathcal{X}}
\newcommand{\cY}{\mathcal{Y}}

\newcommand{\sA}{\mathscr{A}}
\newcommand{\sB}{\mathscr{B}}
\newcommand{\sC}{\mathscr{C}}
\newcommand{\sE}{\mathscr{E}}
\newcommand{\sF}{\mathscr{F}}
\newcommand{\sG}{\mathscr{G}}
\newcommand{\sH}{\mathscr{H}}
\newcommand{\sL}{\mathscr{L}}
\newcommand{\sS}{\mathscr{S}}
\newcommand{\sT}{\mathscr{T}}
\newcommand{\sX}{\mathscr{X}}
\newcommand{\sY}{\mathscr{Y}}
\newcommand{\sZ}{\mathscr{Z}}

\newcommand{\tS}{\tilde{S}}
% Debug
\newcommand{\todo}[1]{\begin{color}{blue}{{\bf~[TODO:~#1]}}\end{color}}

% a few handy macros

\renewcommand{\le}{\leqslant}
\renewcommand{\ge}{\geqslant}
\newcommand\matlab{{\sc matlab}}
\newcommand{\goto}{\rightarrow}
\newcommand{\bigo}{{\mathcal O}}
%\newcommand{\half}{\frac{1}{2}}
%\newcommand\implies{\quad\Longrightarrow\quad}
\newcommand\reals{{{\rm l} \kern -.15em {\rm R} }}
\newcommand\complex{{\raisebox{.043ex}{\rule{0.07em}{1.56ex}} \hskip -.35em {\rm C}}}


% macros for matrices/vectors:

% matrix environment for vectors or matrices where elements are centered
\newenvironment{mat}{\left[\begin{array}{ccccccccccccccc}}{\end{array}\right]}
\newcommand\bcm{\begin{mat}}
\newcommand\ecm{\end{mat}}

% matrix environment for vectors or matrices where elements are right justifvied
\newenvironment{rmat}{\left[\begin{array}{rrrrrrrrrrrrr}}{\end{array}\right]}
\newcommand\brm{\begin{rmat}}
\newcommand\erm{\end{rmat}}

% for left brace and a set of choices
%\newenvironment{choices}{\left\{ \begin{array}{ll}}{\end{array}\right.}
\newcommand\when{&\text{if~}}
\newcommand\otherwise{&\text{otherwise}}
% sample usage:
%\delta_{ij} = \begin{choices} 1 \when i=j, \\ 0 \otherwise \end{choices}


% for labeling and referencing equations:
\newcommand{\eql}{\begin{equation}\label}
\newcommand{\eqn}[1]{(\ref{#1})}
% can then do
%\eql{eqnlabel}
%...
%\end{equation}
% and refer to it as equation \eqn{eqnlabel}.
% some useful macros for finite difference methods:
\newcommand\unp{U^{n+1}}
\newcommand\unm{U^{n-1}}
% for chemical reactions:
\newcommand{\react}[1]{\stackrel{K_{#1}}{\rightarrow}}
\newcommand{\reactb}[2]{\stackrel{K_{#1}}{~\stackrel{\rightleftharpoons}
 {\scriptstyle K_{#2}}}~}
\makeatletter
\def\th@plain{%
\thm@notefont{}% same as heading font
\itshape % body font
}
\def\th@definition{%
\thm@notefont{}% same as heading font
\normalfont % body font
}
\makeatother
\date{}



\title{ Lecture-22: Communicating classes}
\author{}

\begin{document}
\maketitle

%%%%%%%%%%%%%%%%%%%%%%%%%%%%%%
\section{Communicating classes}
%%%%%%%%%%%%%%%%%%%%%%%%%%%%%%
\begin{defn}
For states $x, y \in \sX$, 
it is said that state $y$ is \textbf{accessible} from state $x$ if $p_{xy}^{(n)} > 0$ for some $n \in \Z_+$, and denoted by $x \to y$. 
If two states $x,y \in \sX$ are accessible to each other, they are said to \textbf{communicate} with each other, denoted by $x \leftrightarrow y$. 
\end{defn}
\begin{prop}
Communication is an equivalence relation. 
\end{prop}
\begin{proof} 
Relation on state space $\sX$ is a subset of product of sets $\sX \times \sX$. 
Communication is a relation on state space $\sX$, as it relates two states $x,y \in \sX$. 
To show equivalence, we have to show reflexivity, symmetry, and transitivity of the relation. 
\begin{itemize}
\item [Reflexivity:] Since $P^0 = I$, it follows that $p_{xx}^{(0)} = 1 > 0$ for each state $x\in\sX$, and hence $x\leftrightarrow x$. % and the reflexivity of the relation follows. 
\item [Symmetry:] If $x \leftrightarrow y$, then we know that $x \to y$ and $y \to x$ and hence $y \leftrightarrow x$. 
%Hence, the symmetry of the relation follows. 
%Reflexivity and symmetry are obvious for this relation, 
\item [Transitivity:] For transitivity, suppose $x \leftrightarrow y$ and $y \leftrightarrow z$. 
Let $m,n \in \Z_+$ such that $p_{xy}^{(m)} >0$ and $p_{yz}^{(n)} >0$. 
Then by Chapman Kolmogorov equation, we have
$
p_{xz}^{(m+n)} = \sum_{w \in \sX} p_{xw}^{(m)}p_{wz}^{(n)} \ge p_{xy}^{(m)}p_{yz}^{(n)} > 0. 
$
This implies $x \to z$, and using similar arguments one can show that $z \to x$, and the transitivity follows.  
%\item [Reflexivity:] If this relation has a single element, then it is obvious. 
%If not, then for $x \leftrightarrow y$, we have Since $p_{xy}^{(n)}> 0$ and $p_{yx}^{(m)} > 0$ for some $m, n \in \Z_+$. 
%Therefore, $p_{xx}^{(n+m)} \ge p_{xy}^{(n)}p_{yx}^{(m)} > 0$ and hence we have $x \leftrightarrow x$, implying the reflexivity of the relation.  
\end{itemize}
\end{proof}
\begin{rem}
Since the communication is an equivalence relation, it partitions state space $\sX$ into equivalence classes. 
\end{rem}
\begin{defn}
Each equivalence class of communication relation is called a \textbf{communicating class}. 
A property of states is said to be a \textbf{class property} if for each communicating class $\cC$, 
either all states in $\cC$ have the property, or none do. 
\end{defn}
\begin{defn}
A set of states that communicate are called a \textbf{communicating class}. 
\begin{compactenum}[(i)]
\item A communicating class $\cC$ is called \textbf{closed} if no edges leave this class. 
That is, for all edges $(x,y) \in \cC \times \cC^c$, we have $p_{xy} = 0$. 
\item An \textbf{open} communicating class is not closed, i.e. there exist an edge that leaves this class. 
That is, there exists an edge $(x,y) \in \cC\times \cC^c$ such that $p_{xy} > 0$.  
\end{compactenum}
\end{defn}



%%%%%%%%%%%%%%%%%%%%%%%%%%%%%%
\subsection{Irreducibility and periodicity}
%%%%%%%%%%%%%%%%%%%%%%%%%%%%%%
%A consequence of the previous result is that communicating classes are disjoint or identical. 
\begin{defn} 
A Markov chain with a single class is called an \textbf{irreducible} Markov chain. 
That is, for any two states $x, y \in \sX$, there exists an integer $n \in \Z_+$ such that $p_{xy}^{(n)} > 0$. 
In other words, any state $y$ can be reached from any state $x$ using transitions of positive probability. 
\end{defn}
\begin{defn}
Let $\cT(x) \triangleq \set{n \in \N : p_{xx}^{(n)} > 0}$ be the set of times when the chain can possibly return to the initial state $x$. 
The \textbf{period} of any state $x \in \sX$ is defined as
$
d(x) \triangleq \gcd\cT(x) =  \gcd\{n \in \N : p_{xx}^{(n)} > 0\}.
$
We define $d(x) = \infty$, if $p_{xx}^{(n)} = 0$ for all $n \in \N$. 
A state $x \in \sX$ is called \textbf{aperiodic} if the period $d(x)$ is $1$.  
\end{defn}

\begin{prop}
If $x \leftrightarrow y$, then $d(x) = d(y)$. 
That is, periodicity is a class property.
\end{prop}
\begin{proof}
Consider two states $x\neq y \in \sX$ and integers $m,n \in \N$ be such that $p_{xy}^{(m)}p_{yx}^{(n)} > 0$. 
Suppose $s \in \cT(x)$, that is $p_{xx}^{(s)} > 0$. Then
\meq{2}{ 
&p_{yy}^{(n+m)} \ge p_{yx}^{(n)}p_{xy}^{(m)} > 0,&& p_{yy}^{(n+s+m)} \ge p_{yx}^{(n)}p_{xx}^{(s)}p_{xy}^{(m)} > 0. 
}
Hence $d(y) | n+m$ and $d(y) | n+s+m$, and hence $d(y) | s$ for any $s \in \cT(x)$. 
In particular, it implies that $d(y) | d(x)$. 
By symmetrical arguments, we get $d(x) | d(y)$. 
Hence $d(x) = d(y)$.
%Source of proof is from \emph{Sheldon Ross: Stochastic Processes}.
\end{proof}
\begin{defn}
For an irreducible chain, the period of the chain is defined to be the period which is common to all states. 
An irreducible Markov chain is called \textbf{aperiodic} if the single communicating class is aperiodic. 
\end{defn}
\begin{prop}
If the transition matrix $P$ is aperiodic and irreducible, 
then there is an integer $r_0$ such that $p_{xy}^{(r)}>0$ for all $x,y \in \sX$  and $r \ge r_0$. 
\end{prop}

%%%%%%%%%%%%%%%%%%%%%%%%
\subsection{Transient and recurrent states}
%%%%%%%%%%%%%%%%%%%%%%%%

%We would donate the conditional probability and conditional expectation of a measurable event $A$ starting from state $x$ as 
%\meq{2}{
%&P_x(A) \triangleq P(A | \{X_0 = x\}), && \E_x1_A \triangleq \E[ 1_A | \{X_0 = x\}]. 
%}

%%\begin{defn}
%Let $f_{xy}^{(n)}$ denote the probability that starting from state $x$, the first transition into state $y$ happens at time $n$. Then let
%\begin{align*}f_{xy} = \sum_{n=1}^\infty f_{xy}^{(n)}\end{align*}
%Here $f_{xy}$ would therefore denote the probability of ever entering state $y$ given that we start at state $x$. 
%State $y$ is said to be \textbf{transient} if $f_{yy} < 1$ and \textbf{recurrent} if $f_{yy}=1$. 
%%\end{defn}
%\begin{prop} 
%The total number of visits to a state $y \in \sX$ is denoted by 
%%\eq{
%$N_y = \sum_{n \in \Z_+}1\{X_n = y\}$.
%%}
%Then, for each $m \in \N$ and $i \neq j$, we have 
%\eq{
%P_j\{N_y = m\} &= f_{yy}^{m-1}(1-f_{yy}), ~m \in \N\\
%P_x\{N_y = m\} &=\begin{cases}
%1 - f_{xy} & m = 0,\\
%f_{xy}f_{yy}^{m-1}(1-f_{yy}) & m \in \N.
%\end{cases}
%}
%\end{prop}
%\begin{proof}
%For each $k \in \N$, the time of the $k$th visit to the state $y$ is a stopping time. 
%From strong Markov property, the next return to state $y$ is independent of the past. 
%Hence, each return to state $y$ is an \textit{iid} Bernoulli random variable with probability $f_{yy}$. 
%It follows that the number of visits $y$ is the time for first failure to return. 
%Conditioned on $X_0 = y$, the distribution of $N_y$ is geometric random variable with success probability $1-f_{yy}$. 
%%Hence the probabilities will multiply. 
%%Hence starting from state $y$ we visit the state $y$ with probability $f_{yy}$. 
%%Moreover, we visit the state $y$ exactly $m$ number of times (and no more) with probability $f_{yy}^m(1-f_{yy}).$\\
%
%Conditioned on $X_0 = x$, the stopping time of first visit to $y$ is a Bernoulli random variable with probability $f_{xy}$. 
%Hence, the second result follows. 
%\end{proof}
%
%\begin{cor} 
%For a Markov chain $X$, 
%%\eq{
%$P_j\{N_y < \infty\} = 1\{f_{yy} < 1\}$. 
%%}
%\end{cor}
%\begin{proof}
%We can write the event $\{N_y < \infty\}$ as disjoint union of events $\{N_y = n\}$, to get 
%\eq{
%P_j\{N_y \in \N\} &= \sum_{n \in\N}P_j\{N_y = n\} = 1\{f_{yy} < 1\}.
%}
%\end{proof}
%In particular, this corollary implies the following
%\begin{enumerate}
%	\item A transient state is visited a finite amount of times almost surely.
%	\item A recurrent state is visited infinitely often almost surely.
%	\item In a finite state Markov chain, not all states may be transient. 
%\end{enumerate} 
%\begin{prop}
%A state $y$ is recurrent iff
%%\begin{align*}
%$\sum_{k\in\N} p_{yy}^{(k)} = \infty$.
%%\end{align*}
%\end{prop}
%\begin{proof}
%For any state $y \in \sX$, we can write
%\begin{align*}
%p_{xx}^{(k)} &= \mathbb{P}_i\{X_k=i\} = \E_x 1\{X_k=i\}.
%\end{align*}
%Using monotone convergence theorem to exchange expectation and summation, we obtain  
%%Then, assuming that the infinite summation can be pushed inside the expectation 
%\begin{align*}
%\sum_{k\in \N}p_{xx}^{(k)}=\E_x \sum_{k\in\N}1\{X_k=i\} = \E_x N_i.
%\end{align*}
%Thus, $\sum_{k\in \N} p_{xx}^{(k)}$ represents the expected number of returns $\E_x N_i$ to a state $x$ starting from state $x$, 
%which we know to be finite if the state is transient and infinite if the state is recurrent. 
%%$[\Leftarrow]$ For the reverse argument we will show that if $\sum_{k=1}^{\infty}p_{xx}^{(k)}<\infty$ then the state $x$ will be transient, which should suffice to prove our claim. First we observe that for a transient state $y$, if $M_j$ is the number of visits to the state $y$, then
%%\begin{align*}P[M_j=m|X_0=j]=f_{yy}^m(1-f_{yy}).\end{align*} 
%%
%%Hence now again if we calculate the expected number of times the state $y$ is visited starting from state $y$ can be calculated by the following elementary exercise.
%%\begin{align*}
%%\E_j[M_j=m|X_0=j]= & \sum_{m=1}^{\infty}mf_{yy}^m(1-f_{yy})=\frac{f_{yy}}{(1-f_{yy})}=
%%\begin{cases}
%%\infty \quad if \quad f_{yy}=1 \Leftrightarrow j \quad  recurrent,\\
%%C<\infty \quad if \quad f_{yy} \neq 1 \Leftrightarrow j \quad transient.
%%\end{cases}
%%\end{align*}
%\end{proof}

\begin{prop}
Transience and recurrence are class properties.
\end{prop}
\begin{proof}
Let us start with proving recurrence is a class property. 
Let $x$ be a recurrent state and let $x \leftrightarrow y$. 
Hence there exist some $m,n >0$, such that $P_{xy}^{(m)} > 0$ and $p_{yx}^{(n)}>0$. 
As a consequence of the recurrence, $\sum_{s \in \N} p_{xx}^{(s)} = \infty$. 
$
\sum_{s \in \N} p_{yy}^{(m+n+s)} \ge \sum_{s \in \N} p_{yx}^{(n)} p_{xx}^{(s)} P_{xy}^{(m)} = \infty.
$
%Hence $y$ is recurrent. 
If $x$ were transient instead then $\sum_{s\in\N}p_{xx}^{(s)}<\infty$, and we conclude that $y$ is also transient by the following observation
$
\sum_{s \in \N}  p_{yy}^{(s)} \le \frac{\sum_{s \in \N} p_{xx}^{(m+n+s)} }{ p_{yx}^{(n)} P_{xy}^{(m)}}< \infty.
$
%Hence $y$ is transient.
\end{proof}
\begin{cor}
If $y$ is a recurrent state and there exists a state $x$ such that $y \to x$, then $x \to y$ and $f_{xy} = 1$.
\end{cor}
\begin{proof} 
Let $y \in \sX$ be a recurrent state, and consider state $x \in \sX$ such that $y \to x$. 
We will show that $f_{xy} = 1$ and hence $f_{xy}^{(n)} > 0$ for some $n \in \Z_+$ and $x \to y$. 
To this end, we observe that since $y \to x$, 
there exists an integer $n \in \Z_+$ such that the probability of hitting state $x$ for the first time starting from state $y$ in $n$-steps is positive. 
That is, 
\EQ{
f_{yx}^{(n)} \triangleq P_y\set{X_n = x, X_{n-1} \neq x,\dots, X_{1}\neq x} = P_y\set{\tau_X^{\set{x},1} = n} > 0.
} 
%Suppose $f_{xy} < 1$, then 
From the strong Markov property, we have 
\EQ{
1-f_{yy} = P_y\set{\tau_X^{\set{y},1} = \infty} \ge P_y\set{\tau_X^{\set{y},1} = \infty, \tau_X^{\set{x},1}  = n} = P_x\set{\tau_X^{\set{y},1} = \infty}P_y\set{\tau_X^{\set{x},1}  = n} = f_{yx}^{(n)}(1-f_{xy}). % > 0.
}
%This is a contradiction 
Since state $y$ is recurrent, it follows that $f_{xy} = 1$ and hence $x \to y$. 
\end{proof}
\begin{cor} 
Let $x, y \in \sX$ be in the same communicating class and the state $y$ is recurrent.  
Then, $\lim_{n \in \N}\frac{\sum_{k=1}^np_{xy}^{(k)}}{n} = \frac{1}{\mu_{yy}}$. 
Furthermore, if the state $y$ is aperiodic, then $\lim_{n \in \N}p_{xy}^{(n)} = \frac{1}{\mu_{yy}}$. 
\end{cor}
\begin{proof} 
Since $y$ is recurrent and $y \to x$, it follow that $f_{xy} = 1$ from the previous Lemma. 
From the Theorem 1.7 in previous lecture, it follows that $\lim_{n \in \N}\frac{\sum_{k=1}^np_{xy}^{(k)}}{n} = \frac{1}{\mu_{yy}}$. 

Let the period of the state $y$ be $d$. 
Then we know that there exists a positive integer $r_0$ such that for all $n \ge r_0$, we have $p_{yy}^{(nd)} > 0$. 
%Then, $p_{yy}^{(nd)} = \E_y\indicator{X_{nd} = y} = \E_y\left[N_y(nd) - N_y((n-1)d)\right]$. 
%If the state $y$ is aperiodic, then we can write 
%\EQ{
%p_{xy}^{(n)} = \sum_{k=1}^nP_x\set{H_y = k, X_n = y} = \sum_{k=1}^nf_{xy}^{(k)}p_{yy}^{(n-k)}. 
%}
\end{proof}
\begin{thm} 
The states in a communicating class are of one of the following types; all transient, or all null recurrent, or all positive recurrent. 
\end{thm}
\begin{proof} 
It suffices to show that if $x,y$ belong to the same communicating class and $y$ is null recurrent, then $x$ is null recurrent as well. 
We take $r,s \in \N$, such that $p_{yx}^{(r)}p_{xy}^{(s)} > 0$. 
It follows that $p_{yy}^{r+\ell+s} \ge p_{yx}^{(r)}p_{xx}^{(\ell)}P_{xy}^{(s)}$ for all $\ell \in \N$. 
Hence, for any $n > r + s$, we have 
\EQ{
\frac{1}{n}\sum_{k=1}^{n}p_{yy}^{(k)} \ge \frac{1}{n}\sum_{k=r+s+1}^{n}p_{yy}^{(k)} \ge \left(\frac{n-r-s}{n}\right)\left(\frac{1}{n-r-s}\sum_{\ell = 1}^{n-r-s}p_{xx}^{(\ell)}\right)p_{yx}^{(r)}P_{xy}^{(s)}.
}
Since $y$ is null recurrent LHS goes to zero as $n$ increases, which implies $\lim_{n \in \N}\frac{1}{n}\sum_{\ell = 1}^{n}p_{xx}^{(\ell)} = 0$. 
Hence, $x$ is null recurrent as well. 
\end{proof}

\end{document}
